Most knowledge graph construction tools stop at extracting entities and typed relations: they tell you what a corpus says, but not \emph{how} the corpus says it. A relation between a compound and a disease in a peer-reviewed clinical paper carries different epistemic weight from the same relation in a forward-looking patent claim, yet flat knowledge graphs collapse the distinction. I present Epistract, an open-source domain-agnostic knowledge graph framework distributed as a Claude Code plugin, that introduces a \emph{two-layer} graph: Layer 1 contains the brute facts (entities, typed relations, evidence quotes, confidence scores), and Layer 2 is an \emph{epistemic super-domain layer} that annotates every relation with a categorical status (\texttt{asserted}, \texttt{prophetic}, \texttt{hypothesized}, \texttt{contested}, \texttt{contradiction}, \texttt{negative}, \texttt{speculative}) using a rule engine combining regex patterns over evidence text with document-type inference. A single domain-specific analyst persona, defined as one YAML string, drives two complementary surfaces: a reactive chat interface over the graph, and a proactive auto-generated briefing that synthesizes the classified graph into temporally-stratified, contradiction-aware analyst output. The framework ships with four pre-built domains: drug-discovery (validated across six scenarios spanning neurogenetics, oncology, rare disease, immuno-oncology, cardiovascular inflammation, and GLP-1 competitive intelligence -- 111 documents, 783 nodes, 2{,}230 relations across S1--S5 plus 278 nodes / 855 relations on S6 alone after a v3 rebuild), clinicaltrials (validated on a GLP-1 Phase 3 corpus from \texttt{ClinicalTrials.gov} -- 10 protocols, 142 nodes, 395 relations), contracts, and fda-product-labels (the latter two are schema scaffolds awaiting their first showcase corpus). Each domain carries its own epistemic rules: phase-based evidence grading for clinicaltrials, a four-level FDA epistemology classifier (established / reported / theoretical / asserted) for fda-product-labels. Adding a new domain requires no pipeline changes -- it is four configuration files (schema, extraction prompt, epistemic rules, persona) generated interactively via \texttt{/epistract:domain}. The framework, all corpora, and all artifacts are MIT-licensed and reproducible.\footnote{\url{https://github.com/usathyan/epistract}}
